% !TEX spellcheck = en_US
% !TEX spellcheck = LaTeX
%This should be Lecture-09
\documentclass[letterpaper,10pt,english]{article}
\usepackage{%
amsfonts,%
amsmath,%
amssymb,%
amsthm,%
babel,%
bbm,%
%biblatex,%
caption,%
centernot,%
color,%
enumerate,%
%enumitem,%
epsfig,%
epstopdf,%
etex,%
fancybox,%
framed,%
fullpage,%
%geometry,%
graphicx,%
hyperref,%
latexsym,%
mathptmx,%
mathtools,%
multicol,%
paralist,%
pgf,%
pgfplots,%
pgfplotstable,%
pgfpages,%
proof,%
psfrag,%
%subfigure,%
tcolorbox,%
tfrupee,%
tikz,%
times,%
%ulem,%
url,%
xcolor,%
mathpazo,%
}

\definecolor{shadecolor}{gray}{.95}%{rgb}{1,0,0}
\usepackage[margin=1in,top=0.75in]{geometry}
\usepackage[mathscr]{eucal}
\usepgflibrary{shapes}
\usepgfplotslibrary{fillbetween}
\usetikzlibrary{%
arrows,%
backgrounds,%
chains,%
decorations.pathmorphing,% /pgf/decoration/random steps | erste Graphik
decorations.text,% 
matrix,%
positioning,% wg. " of "
fit,%
patterns,%
petri,%
plotmarks,%
scopes,%
shadows,%
shapes.misc,% wg. rounded rectangle
shapes.arrows,%
shapes.callouts,%
shapes%
}

%\pgfplotsset{compat=newest} %<------ Here
\pgfplotsset{compat=1.11} %<------ Or use this one

\theoremstyle{plain}
\newtheorem{thm}{Theorem}[section]
\newtheorem{lem}[thm]{Lemma}
\newtheorem{prop}[thm]{Proposition}
\newtheorem{cor}[thm]{Corollary}
\newtheorem{clm}[thm]{Claim}

\theoremstyle{definition}
\newtheorem{defn}[thm]{Definition}
\newtheorem{conj}[thm]{Conjecture}
\newtheorem{exmp}[thm]{Example}
\newtheorem{axiom}[thm]{Axiom}
\newtheorem{assum}[thm]{Assumption}
\newtheorem{exerc}[thm]{Exercise}
\newtheorem*{defin}{Definition}

\theoremstyle{remark}
\newtheorem{rem}{Remark}
\newtheorem{note}{Note}

\newcommand{\iid}{\emph{i.i.d.}\ }
\newcommand{\Cov}{\operatorname{Cov}}
%\newcommand{\det}{\operatorname{det}}
%\newcommand{\Real}{\mathbb{R}}
\newcommand{\tr}{\operatorname{tr}}
\newcommand{\Var}{\operatorname{Var}}
\newcommand{\cov}{\operatorname{cov}}

%\renewcommand{\proof}[1]{\begin{proof}#1\end{proof}}
\newcommand{\EQ}[1]{\begin{equation*}#1\end{equation*}}
\newcommand{\EQN}[1]{\begin{equation}#1\end{equation}}
\newcommand{\eq}[1]{\begin{align*}#1\end{align*}}
\newcommand{\meq}[2]{\begin{xalignat*}{#1}#2\end{xalignat*}}
\newcommand{\norm}[1]{\left\lVert#1\right\rVert}
\newcommand{\inner}[1]{\left\langle #1\right\rangle}
\newcommand{\abs}[1]{\left\lvert#1\right\rvert}
\newcommand{\expect}[1]{\mathbb{E}\left[{#1}\right]}
\newcommand{\prob}[1]{\mathbb{P}\left[{#1}\right]}
\newcommand{\given}{\; \big\vert \;} 
\newcommand{\set}[1]{\left\{#1\right\}} 
\newcommand{\indicator}[1]{1_{\set{#1}}} 
\newcommand{\Ind}[1]{\mathbbm{1}_{#1}} 
\newcommand{\SetIn}[1]{\mathbbm{1}_{\set{#1}}} 
\newcommand{\red}[1]{\textcolor{red}{#1}} 
\newcommand{\floor}[1]{\left\lfloor#1\right\rfloor}
\newcommand{\ceil}[1]{\left\lceil#1\right\rceil}


\newcommand{\C}{\mathbb{C}}
\newcommand{\D}{\mathbb{D}}
\newcommand{\E}{\mathbb{E}}
\newcommand{\F}{\mathbb{F}}
\newcommand{\N}{\mathbb{N}}
\renewcommand{\P}{\mathbb{P}}
\newcommand{\Q}{\mathbb{Q}}
\newcommand{\R}{\mathbb{R}}
\newcommand{\Z}{\mathbb{Z}}

\newcommand{\bA}{\mathbf{A}}
\newcommand{\bI}{\mathbf{I}}
\newcommand{\bU}{\mathbf{1}}
\newcommand{\bx}{\mathbf{x}}
\newcommand{\bX}{\mathbf{X}}
\newcommand{\bZ}{\mathbf{Z}}


\newcommand{\cA}{\mathcal{A}}
\newcommand{\cB}{\mathcal{B}}
\newcommand{\cC}{\mathcal{C}}
\newcommand{\cF}{\mathcal{F}}
\newcommand{\cG}{\mathcal{G}}
\newcommand{\cM}{\mathcal{M}}
\newcommand{\cN}{\mathcal{N}}
\newcommand{\cP}{\mathcal{P}}
\newcommand{\cT}{\mathcal{T}}
\newcommand{\cX}{\mathcal{X}}
\newcommand{\cY}{\mathcal{Y}}

\newcommand{\sA}{\mathscr{A}}
\newcommand{\sB}{\mathscr{B}}
\newcommand{\sC}{\mathscr{C}}
\newcommand{\sE}{\mathscr{E}}
\newcommand{\sF}{\mathscr{F}}
\newcommand{\sG}{\mathscr{G}}
\newcommand{\sH}{\mathscr{H}}
\newcommand{\sL}{\mathscr{L}}
\newcommand{\sS}{\mathscr{S}}
\newcommand{\sT}{\mathscr{T}}
\newcommand{\sX}{\mathscr{X}}
\newcommand{\sY}{\mathscr{Y}}
\newcommand{\sZ}{\mathscr{Z}}

\newcommand{\tS}{\tilde{S}}
% Debug
\newcommand{\todo}[1]{\begin{color}{blue}{{\bf~[TODO:~#1]}}\end{color}}

% a few handy macros

\renewcommand{\le}{\leqslant}
\renewcommand{\ge}{\geqslant}
\newcommand\matlab{{\sc matlab}}
\newcommand{\goto}{\rightarrow}
\newcommand{\bigo}{{\mathcal O}}
%\newcommand{\half}{\frac{1}{2}}
%\newcommand\implies{\quad\Longrightarrow\quad}
\newcommand\reals{{{\rm l} \kern -.15em {\rm R} }}
\newcommand\complex{{\raisebox{.043ex}{\rule{0.07em}{1.56ex}} \hskip -.35em {\rm C}}}


% macros for matrices/vectors:

% matrix environment for vectors or matrices where elements are centered
\newenvironment{mat}{\left[\begin{array}{ccccccccccccccc}}{\end{array}\right]}
\newcommand\bcm{\begin{mat}}
\newcommand\ecm{\end{mat}}

% matrix environment for vectors or matrices where elements are right justifvied
\newenvironment{rmat}{\left[\begin{array}{rrrrrrrrrrrrr}}{\end{array}\right]}
\newcommand\brm{\begin{rmat}}
\newcommand\erm{\end{rmat}}

% for left brace and a set of choices
%\newenvironment{choices}{\left\{ \begin{array}{ll}}{\end{array}\right.}
\newcommand\when{&\text{if~}}
\newcommand\otherwise{&\text{otherwise}}
% sample usage:
%\delta_{ij} = \begin{choices} 1 \when i=j, \\ 0 \otherwise \end{choices}


% for labeling and referencing equations:
\newcommand{\eql}{\begin{equation}\label}
\newcommand{\eqn}[1]{(\ref{#1})}
% can then do
%\eql{eqnlabel}
%...
%\end{equation}
% and refer to it as equation \eqn{eqnlabel}.
% some useful macros for finite difference methods:
\newcommand\unp{U^{n+1}}
\newcommand\unm{U^{n-1}}
% for chemical reactions:
\newcommand{\react}[1]{\stackrel{K_{#1}}{\rightarrow}}
\newcommand{\reactb}[2]{\stackrel{K_{#1}}{~\stackrel{\rightleftharpoons}
 {\scriptstyle K_{#2}}}~}
\makeatletter
\def\th@plain{%
\thm@notefont{}% same as heading font
\itshape % body font
}
\def\th@definition{%
\thm@notefont{}% same as heading font
\normalfont % body font
}
\makeatother
\date{}


\title{Lecture-09: Moments and $L^p$ spaces}
%\author{}

\begin{document}
\maketitle

%%%%%%%%%%%%%%%%%%%%%%%%%%%%%%%%%%%%%%%%%%%%%%%%%%%%%
\section{Moments}
%%%%%%%%%%%%%%%%%%%%%%%%%%%%%%%%%%%%%%%%%%%%%%%%%%%%%
Let $X:\Omega\to\R$ be a random variable defined on the probability space $(\Omega,\sF,P)$ with the distribution function $F_X:\R \to [0,1]$. 
\begin{shaded}
\begin{exmp}[Absolute value function] 
For the function $\abs{\cdot}: \R \to \R_+$, 
we can write the inverse image of half open sets $(-\infty, x]$ for any $x \in \R$ as $A_g(x) = g^{-1}(-\infty, x]$. 
It follows that $A_g(x) = \emptyset \in \sB(\R)$ for $x < 0$ and $A_g(x) = [-x, x] \in \sB(\R)$ for $x \in \R_+$. 
Since $g^{-1}(-\infty, x] \in \sB(R)$, it follows that $\abs{\cdot}: \R \to \R_+$ is a Borel measurable function. 
\end{exmp}
\end{shaded}
\begin{lem} 
If $\E\abs{X}$ is finite, then $\E X$ exists and is finite. 
\end{lem}
\begin{proof} 
The function $\abs{\cdot}: \R \to \R$ is a Borel measurable function and hence $\abs{X}$ is a random variable. 
Further $\abs{X} \ge 0$, and hence the expectation $\E\abs{X}$ always exists. 
If $\E\abs{X}$ is finite, it means $\E X_+$ and $\E X_-$ are both finite, 
and hence $\E X = \E X_+ - \E X_-$ is finite as well.  
\end{proof}
\begin{cor} 
Let $g: \R \to \R$ be a Borel measurable function. 
If $\E\abs{g(X)}$ is finite, then $\E g(X)$ exists and is finite. 
\end{cor} 
\begin{tcolorbox}
\begin{exerc}[Polynomial function] 
For any $k \in \N$, we define functions $g_k: \R \to \R$ such that $g_k: x \mapsto x^k$. 
Show that $g_k$ is Borel measurable for all $k \in \N$. 
\end{exerc} 
\end{tcolorbox}
%Let $X: \Omega \to \R$ be a random variable defined on the probability space $(\Omega, \sF, P)$. 
%Recall that if $g: \R \to \R$ be a Borel measurable function, then $g(X)$ is also a random variable. 
\begin{defn}[Moments]
%Let $X: \Omega \to \R$ be a random variable defined on the probability space $(\Omega, \sF, P)$. 
We define the $k$th \textbf{moment} of the random variable $X$ as $m_k \triangleq \E g_k(X) = \E X^k$. 
First moment $\E X$ is called the \textbf{mean} of the random variable. 
\end{defn}
\begin{rem} 
If $\E\abs{X}^k$ is finite, then $m_k$ exists and is finite. 
\end{rem}
\begin{rem}
If $P\set{\abs{X} \le 1} = 1$, then $P\set{\abs{X}^k \le 1} = 1$.  
Therefore, by the monotonicity of expectations $\E\abs{X}^k \le 1$, 
and the moments $m_k$ exist and are finite for all $k \in \N$. 
\end{rem}

%%%%%%%%%%%%%%%%%%%%%%%%%%%%%%%%%%%%
\section{$L^p$ spaces}
%%%%%%%%%%%%%%%%%%%%%%%%%%%%%%%%%%%%
\begin{defn}
A vector space over field $\F$ is denoted by a set $V$ has (a) vector addition $+: V\times V\to V$ that satisfies associativity and commutativity, and has an identity and inverse element for each vector $v\in V$, and (b) scalar multiplication $\cdot: \F \times V \to V$ that satisfies compatibility of field multiplication, distributivity with vector and field addition, and has an identity element. 
\end{defn}
\begin{rem} 
We can verify that the set of random variables is a vector space over reals, and we denote it by $V$. 
To this end, we can verify the associativity and commutativity of addition of random variables. 
For any random variable $X$, the constant function $0$ is the identity random variable, and $-X$ is the additive inverse for vector additions.  
We can verify the compatibility of scalar multiplication with random variables, 
existence of unit scalar $1$, and distributivity of scalar multiplications with respect to field addition and vector addition.  
\end{rem}
\begin{defn} 
For a vector space $V$ over field $\F$, a \textbf{norm} on the vector space is a map $f: V \to \R_+$ 
that satisfies  
\begin{enumerate}
	\item [\textbf{homogeneity:}] $f(av) = \abs{a}f(v)$ for all $a \in \F$ and $v \in V$,
	\item [\textbf{sub-additivity:}] $f(v+w) \le f(v) + f(w)$ for all $v,w \in V$, and 
	\item [\textbf{point-separating:}] $f(v) \ge 0$ for all $v \in V$. 
\end{enumerate}
\end{defn}
\begin{defn}
For a probability space $(\Omega,\sF,P)$, and $p \ge 1$, 
we define the set of random variables with finite absolute $p$th moment as the vector space  
$
L^p \triangleq \set{X \in V:(\E\abs{X}^p)^{\frac{1}{p}} < \infty}. 
$
\end{defn}
\begin{defn}
We define a function $\norm{}_p: L^p \to \R_+$ defined by $\norm{}_p(X) = \norm{X}_p \triangleq (\E\abs{X}^p)^{\frac{1}{p}}$ for any $X \in L^p$ and real $p \ge 1$. 
\end{defn}

\begin{rem}
For $p=1$, the map $\norm{}_p$ is norm. 
Therefore, $L^1$ is a normed vector space consisting of random variables with bounded absolute mean.
\end{rem}
\begin{rem}
We will show that $\norm{X}_\infty = \sup\set{\abs{X(\omega)}:\omega\in \Omega}$, 
and hence $L^\infty$ is a normed vector space of bounded random variables. 
%It follows that the function $\norm{}_p$ is a norm for vector space $L^p$ for $p \in \set{1,\infty}$. 
\end{rem}
\begin{rem}
We will also show that $\norm{}_p$ is a norm for all $p \in (1, \infty)$, 
and hence $L^p$ is a normed vector space of random variables with bounded $\norm{}_p$ norm. 
In particular, the $L^2$ space consists of random variables with bounded second moment.
\end{rem}
\begin{rem}
If $\E\abs{X}^N$ is finite for some $N \in \N$, 
then $\E\abs{X}^k$ is finite for all $k \in [N]$. 
This follows from the linearity and monotonicity of expectations, and the fact that 
\EQ{
\abs{X}^k = \abs{X}^k\SetIn{\abs{X} \le 1} + \abs{X}^k\SetIn{\abs{X} > 1} \le 1 + \abs{X}^N. 
}
This implies that $L^N \subseteq L^k$ for all $k \in [N]$. 
We will show that $L^{q} \subseteq L^p$ for any real numbers $1 \le p \le q$. 
\end{rem}
%\begin{lem} 
%If $m_N$ is finite for some $N \in \N$, then $m_k$ is finite for all $k \in [N]$.
%\end{lem}
%\begin{proof} 
%For any random variable $X: \Omega \to \R$ and $k \in [N]$, we can write 
%\EQ{
%\abs{X}^k = \abs{X}^k\SetIn{\abs{X}^k \le 1} + \abs{X}^k\SetIn{\abs{X}^k > 1} \le \SetIn{\abs{X}^k \le 1} + \abs{X}^N\SetIn{\abs{X}^k > 1} \le 1 + \abs{X}^N. 
%}
%The result follows from the monotonicity of expectations. 
%\end{proof}

%%%%%%%%%%%%%%%%%%%%%%%%%%%%%%%%%%%%
\section{Central Moments}
%%%%%%%%%%%%%%%%%%%%%%%%%%%%%%%%%%%%
\begin{shaded}
\begin{exmp}[Shifted polynomial functions] 
For any $k \in \N$, we define functions $h_k: \R \to \R$ such that $h_k: x \mapsto (x-m_1)^k$. 
Then, $h_k = g_k(x-m_1) = g_k\circ f$ where $f:\R \to\R$ is defined as $f(x) = x-m_1$ for all $x \in \R$. 
Since $g_k$ and $f$ are measurable, so is $h_k$. 
\end{exmp}
\end{shaded}
\begin{defn}[Central moments] 
Let $X: \Omega \to \R$ be a random variable defined on the probability space $(\Omega, \sF, P)$ with finite first moment $m_1$. 
We define the $k$th \textbf{central moment} of the random variable $X$ as $\sigma_k \triangleq \E h_k(X) = \E (X-m_1)^k$. 
%First central moment $\sigma_1= \E(X-m_1) = 0$ and 
The second central moment $\sigma_2 = \E(X-m_1)^2$ is called the \textbf{variance} of the random variable and denoted by $\sigma^2$.  
\end{defn}
\begin{lem} 
The first central moment $\sigma_1= \E(X-m_1) = 0$ and the variance $\sigma^2 = \E(X-m_1)^2$ for a random variable $X$ is always non-negative, 
with equality when $X$ is a constant. 
That is, $m_2 \ge m_1^2$ with equality when $X$ is a constant. 
\end{lem}
\begin{proof}
Recall that $h_1, h_2$ are Boreal measurable functions, and hence $h_1(X) = X- m_1$ and $h_2(X) = (X-m_1)^2$ are random variables. 
From the linearity of expectations, it follows that $\sigma_1 = \E h_1(X) = \E X - m_1 = 0$. 
Since $(X-m_1)^2 \ge 0$ almost surely, it follows from the monotonicity of expectation that $0 \le \E(X-m_1)^2$.  
From the linearity of expectation and expansion of $(X-m_1)^2$, we get 
$
\sigma^2 = \E X^2 - 2m_1\E X + m_1^2 = m_2 - m_1^2 \ge 0.
$
\end{proof}
\begin{rem} 
If second moment is finite, then the first moment is finite. 
That is, $L^2 \subseteq L^1$. 
\end{rem}

%%%%%%%%%%%%%%%%%%%%%%%%%%%%%%%%%%%%
\section{Inequalities}
%%%%%%%%%%%%%%%%%%%%%%%%%%%%%%%%%%%%
\begin{thm}[Markov's inequality]
Let $X: \Omega \to \R$ be a random variable defined on the probability space $(\Omega, \sF, P)$. 
Then, for any monotonically non-decreasing function $f: \R \to \R_+$, we have 
\EQ{
P\set{X \ge \epsilon} \le \frac{\E[f(X)]}{f(\epsilon)}.
}
\end{thm}
\begin{proof} 
We can verify that any monotonically non-decreasing function $f: \R \to \R_+$ is Borel measurable. 
Hence,  $f(X)$ is a random variable for any random variable $X$. 
Therefore, 
\EQ{
f(X) = f(X)\SetIn{X \ge \epsilon} + f(X)\SetIn{X < \epsilon} \ge f(\epsilon)\SetIn{X \ge \epsilon}.
}
The result follows from the monotonicity of expectations. 
\end{proof}
\begin{cor}[Markov] 
Let $X$ be a non-negative random variable, then $P\set{X \ge \epsilon} \le \frac{\E X}{\epsilon}$ for all $\epsilon  > 0$.
\end{cor}
\begin{cor}[Chebychev] 
Let $X$ be a random variable with finite mean $m_1$ and variance $\sigma^2$, 
then 
\EQ{
P\set{\abs{X -m_1} \ge \epsilon} \le \frac{\sigma_2}{\epsilon^2},\text{ for all }\epsilon  > 0.
} 
\end{cor}
\begin{proof} 
Apply the Markov's inequality for random variable $Y = \abs{X-m_1} \ge 0$ and increasing function $f(x) = x^2$ for $x \ge 0$. 
\end{proof}
\begin{cor}[Chernoff] 
Let $X$ be a random variable with finite $\E[e^{\theta X}]$ for some $\theta>0$, 
then 
\EQ{
P\set{X \ge \epsilon} \le e^{-\theta \epsilon}\E[e^{\theta X}],\text{ for all }\epsilon  > 0.
} 
\end{cor}
\begin{proof} 
Apply the Markov's inequality for random variable $X$ and increasing function $f(x) = e^{\theta x} > 0$ for $\theta > 0$. 
\end{proof}
\begin{defn}[Convex function] 
A real-valued function $f: \R \to \R$ is convex if for all $x, y\in \R$ and $\theta \in [0,1]$, we have 
\EQ{
f(\theta x + (1-\theta)y) \le \theta f(x) + (1-\theta)f(y).
}
\end{defn}
\begin{thm}[Jensen's inequality] 
For any convex function $f: \R \to \R$ and random variable $X$, we have 
\EQ{
f(\E X) \le \E f(X).
}
\end{thm}
\begin{proof}
It suffices to show this for simple random variables $X: \Omega \to \sX$. 
We show this by induction on cardinality of alphabet $\sX$.  
The inequality is trivially true for $\abs{\sX}=1$. 
We assume that the inductive hypothesis is true for $\abs{\sX} = n$. 

Let $X \in \sX$, where $\abs{\sX} = n+1$. 
We can denote $\sX = \set{x_1, \dots, x_{n+1}}$ with $p_i \triangleq P\set{X = x_i}$ for all $i \in [n+1]$. 
We observe that $(\frac{p_j}{1-p_1}: j \ge 2)$ is a probability mass function for some random variable $Y \in \sY = \set{x_2, \dots, x_{n+1}}$ with cardinality $n$. 
Hence, by inductive hypothesis, we have 
\EQ{
f\left(\sum_{i=2}^{n+1}\frac{p_i}{1-p_1}x_i\right) 
= f (\E Y) 
\le \E f(Y) 
= \sum_{i=2}^{n+1}\frac{p_i}{1-p_1}f(x_i).
}
%We denote $X = \sum_{i=1}^{n+1}x_i\SetIn{X = x_i}$, where $p_i \triangleq P\set{X = x_i}$.  
%Then, we can form random variables $Y = \sum_{i=2}^{n+1}x_i\SetIn{Y = x_i}$ such that $P\set{Y = x_i} = \frac{p_i}{1-p_1}$ and $Z = X$. 
%Next, we consider a random variable $Z \in \set{x_1, \sum_{i=2}^{n+1}\frac{p_i}{1-p_1}x_i}$ with probability mass function $(p_1, 1-p_1)$. 
Applying the convexity of $f$ to $\theta = p_1, x = x_1, y = \sum_{i=2}^{n+1}\frac{p_i}{1-p_1}x_i$, we get 
\EQ{
f(\E X) 
= f(\sum_{i=1}^{n+1}p_ix_i) 
= f\Big(p_1x_1 + (1-p_1)\sum_{i=2}^{n+1}\frac{p_i}{1-p_1}x_i\Big) 
\le p_1f(x_1) + (1-p_1)f\Big(\sum_{i=2}^{n+1}\frac{p_i}{1-p_1}x_i\Big).
}
From the inductive step, it follows that the RHS is upper bounded by $\E f(X)$, and the result follows.
%From the convexity of $f$ and the inductive step, we can write 
%\EQ{
%f(\E X) 
%= f(\sum_{i=1}^{n+1}p_ix_i) 
%= f\left(p_1x_1 + (1-p_1)\sum_{i=2}^{n+1}\frac{p_i}{1-p_1}x_i\right) 
%= f(\E Z)
%\le \E f(Z) 
%= \sum_{i=1}^{n+1}p_if(x_i) 
%= \E f(X).
%}
\end{proof}
\begin{thm} For any real numbers $1 \le p \le q < \infty$, we have $L^q \subseteq L^p$.
\end{thm}
\begin{proof}
Let $1 \le p \le q < \infty$ and consider a convex function $g: \R_+ \to \R_+$ defined by $g(x) \triangleq x^{{q}/{p}}$ for all $x \in \R_+$. 
It follows that $g(\abs{X}^p) = \abs{X}^q$ and hence from the Jensen's inequality, we get
\EQ{
\norm{X}_p^{q} = g(\E\abs{X}^p) \le \E g(\abs{X}^p) = \norm{X}_q^q.
} 
\end{proof}
\end{document}
